\setOLPbooktitle[সেট, যুক্তি,\\ গণনা]{সেট, যুক্তি, গণনা}
\setOLPbooksubtitle{অধিযুক্তিবিদ্যার একটি উন্মুক্ত ভূমিকা}
\setOLPauthor{Richard Zach}
\setOLPbookversion{ভারতীয় বাংলা সংস্করণ}{বাংলা (ভারত)}
\setOLPsourcelink{https://slc.openlogicproject.org/}
\setOLPauthorlink{https://richardzach.org/}
\AtBeginDocument{%
  \renewcommand*{\olremixedby}{%
    মূল পাঠের রচয়িতা ও সংকলক: \href{\OLPauthorlink}{\OLPauthor}\par
    ওপেন লজিক প্রকল্পের পাঠের উপর ভিত্তি করে।\par\medskip
    ভারতীয় বাংলা পাঠ: OpenAI Codex-এর এআই অনুবাদ।\par
    এই পাঠ্যবইয়ের অনুবাদ, সম্পাদনা ও অক্ষরবিন্যাস:
    GPT-6.1 Sol (Ultra)।\par
    পুনর্ব্যবহৃত বাংলা ওপেন লজিক পাঠে GPT-5.6 Sol ও GPT-6 Sol-এর
    অনুবাদ এবং GPT-6.1 Sol-এর সম্পাদনা রয়েছে (Ultra)।}%
  \renewcommand*{\olcolophon}{বাংলা অক্ষর Noto Serif Bengali-তে এবং
    লাতিন অক্ষর Latin Modern-এ; গাণিতিক অক্ষর মূল বইয়ের রীতি অনুসারে।
    মুদ্রাক্ষরবিন্যাস XeLaTeX-এ।}%
  \renewcommand*{\olbookrevision}{মূল পাঠ: \href{https://github.com/rzach/sets-logic-computation/tree/6178f32f61645f3b62ef621b4c5a88f8e67778ad}{Richard Zach, \emph{Sets, Logic, Computation}}।
    ভিত্তি: \href{https://github.com/OpenLogicProject/OpenLogic/tree/9620cc73f9c8e0ad003c514a5d3748f29611c4c0}{\emph{Open Logic Text}}।
    পাঠের লাইসেন্স CC BY 4.0; প্রচ্ছদের প্রতিকৃতি ও জীবনীতে ব্যবহৃত
    আলোকচিত্রের স্বত্ব ও ব্যবহারের শর্ত পৃথক। আলোকচিত্রের স্বীকৃতি
    বইয়ের শেষে।}%
  \renewcommand*{\oluselicense}[2]{%
    \parbox[b]{.67\linewidth}{\footnotesize
      এই বইয়ের পাঠ \href{https://creativecommons.org/licenses/by/4.0/}{CC BY 4.0}
      লাইসেন্সে প্রকাশিত। এটি Richard Zach-এর \emph{Sets, Logic, Computation}
      ও ওপেন লজিক প্রকল্পের পাঠের উপর ভিত্তি করে।}\hfill
    \href{https://openlogicproject.org/}{%
      \includegraphics[width=.28\linewidth]{\olpath/assets/logos/openlogic-logo-bw}}}%
}
